\chapter{Catatan Asumsi dan Temuan}
\label{ch:09-catatan-asumsi}

\section{Empat Temuan Sisa}

Bab ini mencatat empat hal yang ditemukan saat rancangan diturunkan dari kode yang
benar-benar berjalan. Keempatnya bukan kegagalan yang menghentikan sistem: sistem
berjalan dan seluruh kebutuhan pada \chapref{ch:08-end-to-end} terpenuhi. Yang
dicatat adalah ketidaksesuaian antara apa yang tampak pada skema atau dokumentasi
dan apa yang sesungguhnya terjadi saat runtime.

\subsection{Kolom Ambang Pakan Tidak Dipakai Perangkat}

\textbf{Gejala.} Ambang pakan menipis disimpan pada kolom
\code{low_feed_threshold_pct}, dikirim ke perangkat lewat \code{set_threshold}, dan
dilaporkan kembali di dalam snapshot \code{get_state}. Namun perangkat tidak pernah
memakai nilai itu untuk keputusan apa pun; satu-satunya penjagaan level yang dimiliki
perangkat adalah pemeriksaan apakah persentase sama dengan nol.

\textbf{Dampak.} Mengubah ambang hanya mengubah kapan platform membuat notifikasi
pakan menipis. Hopper yang tersisa lima persen tetap menjalankan jadwal lokal
meskipun ambang disetel dua puluh persen, dan nilai ambang pada snapshot perangkat
dapat disalahartikan sebagai ambang yang berlaku di perangkat.

\textbf{Rekomendasi.} Pilih satu arah yang konsisten: bila ambang ingin ditegakkan di
perangkat, tambahkan penjagaan pada jalur jadwal dan pakan manual sehingga level di
bawah ambang melewati pemberian pakan dan dilaporkan sebagai event; bila ambang
hanya milik platform, hilangkan nilai itu dari snapshot agar tidak menimbulkan kesan
ikut menentukan.

\subsection{Kolom Penaut Riwayat Terjadwal Selalu Kosong}

\textbf{Gejala.} Tabel \code{feeding_logs} memiliki kolom \code{schedule_id}
berkunci asing ke \code{schedules} dengan perilaku \code{SET NULL}, tetapi jalur
penulisan riwayat terjadwal selalu menuliskan null karena event jadwal dari
perangkat tidak membawa pengenal baris jadwal.

\textbf{Dampak.} Riwayat terjadwal tidak dapat ditautkan kembali ke jadwal yang
menghasilkannya; halaman History tidak dapat menampilkan nama jadwal untuk baris
terjadwal. Perilaku \code{SET NULL} yang dirancang untuk menjaga riwayat ketika
jadwal dihapus tidak pernah terpakai, sehingga sebagian skema tidak pernah diuji
data nyata.

\textbf{Rekomendasi.} Sertakan pengenal jadwal pada event \code{scheduled_feed},
karena entri jadwal di perangkat sudah menyimpan nilai itu, lalu isi
\code{schedule_id} saat handler menulis baris riwayat. Bila penautan tidak
diinginkan, hapus kolom beserta kunci asingnya agar skema tidak menjanjikan sesuatu
yang tidak pernah diisi.

\subsection{Status Awal Perintah Tidak Pernah Ditulis}

\textbf{Gejala.} Batasan nilai pada kolom status perintah mengizinkan lima nilai
termasuk \code{pending}, dan sapuan timeout memeriksa \code{pending} maupun
\code{sent}, tetapi fungsi penyisipan selalu menuliskan status terkirim bahkan
sebelum pesan benar-benar dipublikasikan.

\textbf{Dampak.} Nilai \code{pending} tidak pernah muncul pada data nyata. Yang lebih
penting, tidak ada catatan bahwa sebuah perintah sudah tersimpan tetapi belum keluar
dari platform; bila publikasi gagal, baris tetap berstatus terkirim sampai batas
waktu habis.

\textbf{Rekomendasi.} Tulis baris dengan status \code{pending}, publikasikan pesan,
lalu ubah menjadi \code{sent} setelah publikasi terkonfirmasi, sehingga kegagalan
publikasi dapat dibedakan dari perangkat yang tidak menjawab. Alternatif yang lebih
sederhana adalah menghapus nilai itu dari batasan sekaligus dari sapuan.

\subsection{Ambang Offline Ada karena Perilaku Broker Tertanam}

\textbf{Gejala.} Platform menandai perangkat offline setelah \OfflineS{} detik tanpa
data, dan konfigurasi menyebutkan alasan eksplisitnya: broker tertanam tidak menutup
sesi client yang diam. Tanpa penanda itu, perangkat yang mati tanpa memutus soket
akan tetap terlihat online selamanya.

\textbf{Dampak.} Status online pada dashboard adalah kesimpulan platform tentang
kesegaran data, bukan pernyataan broker tentang keadaan sesi. Perangkat yang hidup
tetapi berhenti mengirim telemetri akan tampak offline, dan LWT yang datang
belakangan tidak menambah informasi baru karena statusnya sudah ditetapkan lebih
dahulu.

\textbf{Rekomendasi.} Aktifkan pemeriksaan keepalive pada broker tertanam sehingga
sesi yang diam benar-benar ditutup dan LWT diterbitkan tepat waktu, lalu longgarkan
ambang penanda sebagai lapis kedua. Bila broker tidak dapat diubah, tegaskan pada
dokumentasi bahwa kolom status berarti ``data terakhir diterima kurang dari
\OfflineS{} detik lalu'', bukan ``sesi MQTT masih terbuka''.

\section{Asumsi Lingkungan}

\tabref{tab:09-asumsi} memuat asumsi lingkungan yang menjadi dasar rancangan,
beserta akibatnya bila tidak berlaku. Asumsi yang tidak berlaku biasanya muncul
sebagai gejala yang membingungkan, bukan sebagai pesan galat yang jelas.

\begin{table}[H]
\centering
\caption{Asumsi lingkungan dan akibatnya bila tidak berlaku.}
\label{tab:09-asumsi}
\small
\begin{tabular}{@{}p{3.9cm}p{4.9cm}p{6.2cm}@{}}
\toprule
\textbf{Asumsi} & \textbf{Dasar} & \textbf{Akibat bila tidak berlaku} \\
\midrule
Perangkat dan host berada pada satu jaringan lokal & Alamat broker diisi langsung
sebagai alamat IP host & Perangkat tidak dapat menyambung; alamat harus diperbarui
setiap kali berubah. \\
Identitas perangkat diturunkan dari MAC & Tiga byte terakhir MAC dipakai sebagai
akhiran identitas & Papan pengganti menjadi perangkat baru dengan riwayat kosong;
baris lama tetap ada dan tampak offline. \\
Satu perangkat per sistem & Baris awal tunggal dan pengujian satu node & Perangkat
kedua muncul sebagai baris baru; perilaku banyak perangkat belum diuji. \\
Sistem berkas dapat hilang & Basis data satu berkas di direktori aplikasi & Pada host
berberkas sementara, basis data dibuat ulang dan diisi data awal setiap kali proses
dijalankan; riwayat hanya bertahan bila diunduh lebih dahulu. \\
Waktu dijaga perangkat sendiri & RTC dengan cadangan sel koin dan sinkronisasi NTP &
Tanpa RTC dan NTP, jadwal lokal tidak berjalan sampai jam disetel dari dashboard. \\
Broker tidak memerlukan autentikasi & Broker tertanam menerima koneksi anonim & Siapa
pun di jaringan yang sama dapat mengirim perintah pakan. \\
\bottomrule
\end{tabular}
\end{table}

\section{Batasan Keamanan}

Rancangan ini tidak memiliki kendali keamanan pada lapisan pesan maupun antarmuka.
Yang dimaksud bukan hanya ketiadaan TLS, melainkan ketiadaan seluruh mekanisme
berikut.
\begin{itemize}[leftmargin=*,itemsep=2pt]
  \item \textbf{Tanpa autentikasi broker.} Broker menerima koneksi tanpa nama
        pengguna dan kata sandi. Prefix topik memisahkan proyek, tetapi bukan
        kontrol akses: siapa pun yang menebak prefix dan identitas perangkat dapat
        berlangganan telemetri dan mempublikasikan perintah.
  \item \textbf{Tanpa enkripsi.} Seluruh payload, termasuk pengenal perangkat dan
        jadwal pakan, melintas sebagai teks pada jaringan.
  \item \textbf{Tanpa autentikasi dashboard.} Tidak ada halaman masuk dan tidak ada
        pembedaan pengguna; siapa pun yang menjangkau alamatnya dapat mengirim
        perintah dan mengubah jadwal.
  \item \textbf{Tanpa otorisasi per perintah.} Perangkat menjalankan perintah apa pun
        yang berbentuk benar pada topiknya, tanpa memeriksa asal pengirim.
  \item \textbf{Tanpa idempotensi.} Perintah yang tiba dua kali akan dijalankan dua
        kali bila gate sudah kembali idle.
  \item \textbf{Kredensial di luar repositori.} Nama jaringan dan alamat broker
        disimpan pada berkas terpisah yang tidak ikut disimpan, tetapi berkas itu
        tetap berupa teks biasa pada mesin pengembang.
\end{itemize}

Kesimpulan praktisnya: sistem layak dipakai pada satu jaringan lokal yang seluruh
penghuninya dipercaya, dan tidak layak dipaparkan ke jaringan publik tanpa
autentikasi, TLS, serta lapisan otorisasi di depan topik perintah. Batasan ini
disengaja dan konsisten dengan ruang lingkup pada \chapref{ch:01-pendahuluan}.

\section{Yang Tidak Diklaim Dokumen Ini}

Tiga hal sengaja tidak diklaim. Pertama, dukungan banyak perangkat: skema dan
protokolnya tidak menghalangi, tetapi belum diuji. Kedua, angka latensi round-trip
pada jaringan lokal: yang dijamin hanyalah batas atas \CmdTimeoutS{} detik untuk
setiap perintah. Ketiga, ketahanan terhadap perintah ganda: penjagaan kesibukan gate
mengurangi kemungkinan, tetapi tidak menghilangkan akibat perintah yang tiba dua kali
pada waktu berbeda.
